% Generated draft: requires mathematical and source review.
\documentclass[11pt]{article}
\usepackage[margin=25mm]{geometry}
\usepackage{fontspec}
\setmainfont{Latin Modern Roman}
\usepackage{amsmath,amssymb,mathtools,unicode-math}
\setmathfont{Latin Modern Math}
\usepackage{enumitem,needspace,xurl,hyperref}
\setlength{\parindent}{0pt}
\setlength{\parskip}{5pt}
\setlength{\emergencystretch}{3em}
\newcommand{\N}{\mathbb{N}}
\newcommand{\Z}{\mathbb{Z}}
\newcommand{\Q}{\mathbb{Q}}
\newcommand{\R}{\mathbb{R}}
\newcommand{\C}{\mathbb{C}}
\newcommand{\F}{\mathbb{F}}
\newcommand{\A}{\mathbb{A}}
\newcommand{\PP}{\mathbb{P}}
\newcommand{\E}{\mathbb{E}}
\newcommand{\eps}{\varepsilon}
\newcommand{\dd}{\,\mathrm{d}}
\begin{document}
\section*{UnsolvedMath ID: 20002486}
\textbf{Source problem code:} catalogue entry
\subsection{Definitions and mathematical statement}
Let \((M^{2n},J)\) be an almost complex manifold, \(J^2=-I\). A pseudoholomorphic line in the source is a map \(f:\R^2\to M\) satisfying
\[
 f_y=J(f)f_x.
\]
\textbf{Problem 30 (Victor Bangert).} Develop a global, nonperturbative theory of foliations or laminations by such curves. The source attributes to Moser a perturbative theorem on \(\mathbb T^{2n}\): when \(J\) is close to a standard complex structure and \(v\in\R^{2n}\) is sufficiently irrational, there is a pseudoholomorphic foliation conjugate to a linear foliation by real two-planes containing \(v\). The topology/smallness and arithmetic hypotheses, regularity of conjugacy and leaf parametrizations are not supplied. The source emphasizes \(\mathbb T^4\), comparing positivity of pseudoholomorphic intersection with Hedlund's intersection-based theory of minimal geodesics on \(\mathbb T^2\). It does not state a universal existence theorem for arbitrary \(J\).

\textbf{Reported Chern obstruction in dimension four.} Let \(P\subseteq\R^4\) be a constant real two-plane and let a \(C^1\) diffeomorphism \(F:\mathbb T^4\to\mathbb T^4\) send the linear plane foliation to a foliation whose tangent planes are \(J\)-invariant. The aid reports the necessary condition
\[
 c_1(T\mathbb T^4,J)=0\in H^2(\mathbb T^4;\Z).
\]
This uses the global transported tangent and normal framings of a linearly conjugate foliation; no weaker merely topological conjugacy is substituted.

For a reported family violating this condition, trivialize \(T\mathbb T^4\cong\mathbb T^4\times\R^4\) and identify \(\R^4\) with the quaternions. The orthogonal complex structures given by multiplication by unit imaginary quaternions form a twistor sphere \(S^2\). Pull this family back by
\[
 \mathbb T^4\xrightarrow{\pi}\mathbb T^2\xrightarrow{g}S^2,
 \qquad\deg(g)\ne0.
\]
The aid reports that the resulting smooth almost complex structures have nonzero first Chern class, and therefore admit no foliation of the specified \(C^1\)-conjugate linear type. Nonzero degrees give an infinite obstruction family. This does not exclude every possible pseudoholomorphic foliation or lamination lacking such a conjugacy.

\textbf{Reported slope condition under taming.} If a closed two-form \(\omega\) tames \(J\), meaning \(\omega(v,Jv)>0\) for every nonzero tangent vector, orient the foliation by \(J\). The linear foliation's translation-invariant transverse measure transports through \(F\), defining a Ruelle--Sullivan homology class \(h\in H_2(\mathbb T^4;\R)\). The aid reports
\[
 \langle[\omega],h\rangle>0.
\]
The choice of orientation and transported transverse measure is part of this slope statement; no positivity is asserted without the taming hypothesis.

\textbf{Reported local graph equations.} In a transverse coordinate chart \((x_1,x_2,y_1,y_2)\), write
\[
 J=\begin{pmatrix}A&B\\C&D\end{pmatrix},\qquad
 X_i=\partial_{x_i}+\sum_{a=1}^2P_{ai}(x,y)\partial_{y_a}.
\]
For a \(C^1\) graph-plane field, \(J\)-invariance is the algebraic Riccati equation
\[
 C+DP-PA-PBP=0.
\]
Its integrability is the commuting-frame condition
\[
 \partial_{x_i}P_{aj}+\sum_bP_{bi}\partial_{y_b}P_{aj}
 -\partial_{x_j}P_{ai}-\sum_bP_{bj}\partial_{y_b}P_{ai}=0
 \quad(a=1,2;\ i,j=1,2).
\]
These equations concern graph charts where transversality holds. They do not supply global chart patching, irrationality control or a nonperturbative existence/classification theorem. All supplementary obstructions and equations are reported aid results requiring independent verification.
\subsection{Short English statement}
Bangert asks for a global theory beyond Moser's perturbative foliations. The aid supplies necessary Chern and tame-slope conditions for a specific \(C^1\)-conjugate linear class on \(\mathbb T^4\), plus local graph equations. These restrict that class without ruling out all other foliations or solving the global program.
\subsection{Sources}
\begin{itemize}
\item Local UnsolvedMath registry, permanent ID \(20002486\), catalogue entry, \texttt{aim-pdes-notes.json}, \texttt{notes[97]}, \emph{New connections between dynamical systems and PDE's}, Problem 30, Victor Bangert: \url{https://aimath.org/WWN/dynpde/dynpde.pdf}.
\item Original attributions: Moser's perturbative pseudoholomorphic foliations and Hedlund's minimal-geodesic theory. No full bibliographic references or exact perturbative hypotheses are supplied.
\item Supplied research summary on the Chern obstruction, pulled-back twistor family, transported Ruelle--Sullivan slope and graph equations. Attempt 1: partial\_result, status at run ``open'', rigor ``proved'', low novelty confidence; registry status partially\_solved. The machine-generated aid is not peer reviewed and requires independent verification before citation. The supplied verification identifies a lost leading digit and page artifacts; this draft has not visited the PDF.
\end{itemize}
\end{document}
